Self-critique is increasingly used to improve language model reasoning, yet it remains unknown whether this technique induces genuine restructuring of internal representations or merely produces surface-level variation. We investigate this question by extracting residual stream activations from \pythia across all 33 layers while the model processes 150 \gsm math problems under three matched conditions: base (question and answer), critique (base plus self-critique text), and paraphrase control (base plus length-matched paraphrase). We find that self-critique induces significantly larger representational drift than paraphrasing---critique activations show 41\% greater $L_2$ displacement and lower cosine similarity to base activations across 29 of 33 layers (Bonferroni-corrected $p < 0.05$). The drift is concentrated in middle layers (12--16) and is more directionally consistent across samples (mean cosine consistency 0.596 vs.\ 0.479 for paraphrase), suggesting a systematic ``critique subspace'' rather than random perturbation. Population-level analysis via Centered Kernel Alignment confirms that critique creates genuinely different representational structure (CKA = 0.224 vs.\ 0.428 for paraphrase). These results demonstrate that self-critique produces measurable geometric restructuring of the residual stream, supporting its use in recursive self-improvement architectures and providing a quantitative internal correlate of reflective processing in language models.
